\section{Results}\label{sec:results}

\subsection{The same design-level fault is not the same physical fault}\label{sec:res-severity}

For the same tap pair the SCIG's fault current is about 10\,\% lower than the PMSG's (medians 3.3 against 3.7\,A for inter-turn and 7.7 against 8.5\,A for inter-winding pairs; Online Resource~1, Fig.~S1), but the difference that matters for detection is relative to what the stator sensors see: the fault current is 1.45 times the pre-fault stator current (RMS to RMS, median over pairs) in the PMSG and 0.63 times in the SCIG for inter-turn faults, and 3.3 against 1.5 times for inter-winding faults, because the SCIG's stator current carries the magnetising component (5.3 against 2.6\,A RMS before the fault). The WFSG trials at the matching impedance give 1.20 (inter-turn, 11.6--12.1\,\%) and 1.03 (inter-winding, 2.3--10.6\,\%), and 0.65 at the lower impedance used for its small inter-turn extents. With a series resistance of the order of the winding impedance these currents are resistor-limited: essentially the voltage induced in the shorted turns divided by the resistance, so they scale with volts per turn and speed and differ between the machines mainly through their terminal voltage per turn. A bolted fault would be limited by the leakage reactance of the shorted turns, reach several times rated current and---the first-order difference between the topologies---could be de-excited in the SCIG and the WFSG but not in the PMSG. A fault specified as ``$e$\,\% of a winding'' therefore arrives at the sensors of the three alternatives as three different physical events, and the benchmark's emulation fixes which; Section~\ref{sec:res-mde} uses both axes.

\subsection{Where the signature appears, and why the SDR differs}\label{sec:res-where}

Table~S10 of Online Resource~1 gives the null of the SDR for the main features. It differs by up to a factor of thirty between the PMSG and the SCIG, and not only in sensed quantities: the null of the duty-cycle negative sequence, a controller output that involves no sensor, is five times larger on the PMSG, so the healthy non-stationarity is a property of the machine--drive combination, not of the transducers alone. Table~\ref{tab:decomp} decomposes the SCIG/PMSG ratio of median SDR: for inter-turn faults the PMSG's relative changes are the \emph{larger} ones on every feature (SCIG/PMSG ratio of median $|\delta|$ 0.37--0.71), and the SCIG's higher SDR is entirely a quieter null (ratio of null quantiles 2.6--35); for inter-winding faults the numerators are comparable (0.3--1.4) and the denominators again decide. The SDR is meant to combine both, but the comparison of alternatives must be read with the decomposition beside it.

\begin{table}[t]
\scriptsize
\setlength{\tabcolsep}{3pt}
\caption{Decomposition of the SCIG/PMSG ratio of median SDR into its numerator (ratio of median $|\delta|$, the relative change under fault) and its denominator (ratio $q_{0.95}^{\mathrm{PMSG}}/q_{0.95}^{\mathrm{SCIG}}$ of the null quantiles); five-cycle windows.}\label{tab:decomp}
\begin{tabular}{@{}llccccccc@{}}
\toprule
Class & Feature & $|\delta|$ P & $|\delta|$ S & $q_{95}$ P & $q_{95}$ S & Num. & Den. & SDR ratio \\
\midrule
inter-turn & $V_2/V_1$ & 0.257 & 0.136 & 0.086 & 0.033 & 0.53 & 2.64 & 1.40 \\
inter-turn & $I_2/I_1$ & 0.238 & 0.152 & 0.162 & 0.034 & 0.64 & 4.82 & 3.08 \\
inter-turn & $I_d$\,2$f_e$ & 0.291 & 0.107 & 0.563 & 0.016 & 0.37 & 35.09 & 12.95 \\
inter-turn & PI$_d$\,2$f_e$ & 0.289 & 0.107 & 0.557 & 0.016 & 0.37 & 34.71 & 12.91 \\
inter-turn & $V_q$\,2$f_e$ & 0.222 & 0.155 & 0.165 & 0.060 & 0.70 & 2.73 & 1.91 \\
inter-turn & $V_q^{\mathrm{cmd}}$\,2$f_e$ & 0.197 & 0.139 & 0.162 & 0.039 & 0.71 & 4.17 & 2.94 \\
inter-turn & $D_2/D_1$ & 0.251 & 0.111 & 0.090 & 0.010 & 0.44 & 8.96 & 3.98 \\
inter-winding & $V_2/V_1$ & 5.318 & 1.569 & 0.086 & 0.033 & 0.30 & 2.64 & 0.78 \\
inter-winding & $I_2/I_1$ & 1.833 & 1.869 & 0.162 & 0.034 & 1.02 & 4.82 & 4.91 \\
inter-winding & $I_d$\,2$f_e$ & 2.335 & 1.188 & 0.563 & 0.016 & 0.51 & 35.09 & 17.85 \\
inter-winding & PI$_d$\,2$f_e$ & 2.344 & 1.188 & 0.557 & 0.016 & 0.51 & 34.71 & 17.60 \\
inter-winding & $V_q$\,2$f_e$ & 3.688 & 1.495 & 0.165 & 0.060 & 0.41 & 2.73 & 1.11 \\
inter-winding & $V_q^{\mathrm{cmd}}$\,2$f_e$ & 0.845 & 1.184 & 0.162 & 0.039 & 1.40 & 4.17 & 5.84 \\
inter-winding & $D_2/D_1$ & 2.942 & 0.854 & 0.090 & 0.010 & 0.29 & 8.96 & 2.60 \\
\botrule
\end{tabular}
\end{table}


Fig.~\ref{fig:features} shows the median SDR of every feature for the three alternatives, and Table~S12 of Online Resource~1 the share of trials in which each measurement group carries the strongest signature. In the PMSG the terminal-voltage group is strongest in 67\,\% of inter-turn and 88\,\% of inter-winding trials, through $V_2/V_1$ (median SDR 3.0 and 62 at five-cycle windows; 76 and 94\,\% of trials detectable), and the stator-current negative sequence $I_2/I_1$ is weak for inter-turn faults (median SDR 1.5, 56\,\%). In the SCIG the controller group is strongest in 56 and 67\,\% of trials, through the 2$f_e$ ripple of the $d$-axis current and its PI output (6.7 and 74; 79 and 100\,\%), the stator-current group in 28--33\,\% and the terminal-voltage group in 13\,\% and none. The two machines rank the 30 common features with a Spearman correlation of 0.70 (inter-turn) and 0.89 (inter-winding): the same features are informative in both, in a different order. In the WFSG the strongest inter-turn features are $V_2/V_1$, the $q$-axis current ripple and the commanded $q$-axis voltage, and for inter-winding faults the $q$-axis ripple, the commanded voltage, the duty cycles and the third current harmonic; its field-current 2$f_e$ ripple, an observable the other two machines lack, detected 86\,\% of the inter-turn trials at the matching impedance and all inter-winding trials (Online Resource~1, Table~S1). The physical quantity behind the relocation is the incremental negative-sequence impedance $\Delta V_2/\Delta I_2$ seen at the terminals under inter-winding faults: 15\,$\Omega$ in the PMSG, 7\,$\Omega$ in the SCIG and 11\,$\Omega$ in the WFSG (Online Resource~1, Table~S7), i.e.\ the PMSG drive converts a given negative-sequence current disturbance into twice the terminal voltage that the SCIG drive does; for inter-turn faults the PMSG's negative-sequence current increments are of mixed sign with a median indistinguishable from zero, which is why current-side features see little there. The two machines run the same field-oriented controller; what differs is the angle source (encoder against flux observer), the machine's own negative-sequence impedance and the baseline of each feature, so the relocation is a property of the machine--drive combination rather than of the control structure alone.

\begin{figure}[t]
\centering
\includegraphics[width=\textwidth]{G_sdr_by_feature_3alt}
\caption{Median signal-to-drift ratio of every feature for the three alternatives (dot) with interquartile range (bar), three-cycle windows; features to the right of the dashed line are detectable in more than half of the trials. The WFSG rows pool all its fault impedances.}\label{fig:features}
\end{figure}

\subsection{Minimum detectable extent}\label{sec:res-mde}

Table~\ref{tab:mde} gives the minimum detectable extent with its bootstrap interval, the detectable share with its Wilson interval and the median SDR with its tap-pair bootstrap interval for four representative features; Fig.~\ref{fig:severity} plots the SDR against the design-level extent (left) and the physical severity axis (right). Through the terminal-voltage negative sequence the PMSG detects inter-turn faults from 3.8\,\% [2.8, 3.8] and the SCIG from 7.4\,\% [7.4, 7.4]; through the stator-current negative sequence both need 7.4\,\%, and through the $d$-axis ripple the PMSG needs 12.1\,\% and the SCIG 7.4\,\% [3.8, 7.4]. The smallest tested inter-turn level (2.7--2.8\,\%) is not reliably detectable on either machine by any single feature: the SCIG's $V_2/V_1$ and PI ripple exceed one at 2.7\,\% but not at 2.8 and 3.8\,\%, which the limit-of-detection convention does not credit. For inter-winding faults the SCIG detects from the smallest tested extent (2.3\,\%) on every feature except the $q$-axis voltage ripple, while the PMSG needs 2.3\,\% [2.3, 14.3] on $V_2/V_1$ but 9.7\,\% on $I_2/I_1$ and 21.7\,\% on the $d$-axis ripple. The SDR rises with extent in every machine (Spearman 0.51--0.95), and the MDE intervals are wide because the tested extents form three clusters (about 3, 7.5 and 12\,\% for inter-turn faults), so the SCIG/PMSG ratio of median SDR is the better basis for comparison: 3.1 [1.3, 5.6] for $I_2/I_1$ and 12 [2.6, 30] for the $d$-axis ripple on inter-turn faults, but 1.1 [0.4, 4.6] for $V_2/V_1$ and 1.5 [0.5, 3.2] for the $q$-axis voltage ripple. The two machines differ where the signature relocates and are indistinguishable where it does not. On the physical axis (Fig.~\ref{fig:severity}, right) the picture is the same and independent of the tap-pair design: at equal fault-to-stator current ratio the SCIG's current-side features lie above the WFSG's, which lie above the PMSG's, while on $V_2/V_1$ the three machines nearly coincide. For the WFSG, whose small inter-turn extents were tested only at the lower impedance, the design-level MDE at the matching impedance is bounded by the smallest extent tested there (11.6\,\%, detectable in 78--98\,\% of trials depending on the suite); its 2.7--7.7\,\% inter-turn trials at 1\,$\Omega$ sit at SDR $\approx1$ on $V_2/V_1$ and $I_2/I_1$, and all its inter-winding extents from 2.3\,\% are detectable.

\begin{table}[t]
\footnotesize
\caption{Minimum detectable extent (MDE, \% of the winding between the taps, bootstrap 95\,\% interval in brackets) and detectable share (DS, \%) of four representative features, PMSG and SCIG, five-cycle windows, $\alpha=0.95$. The smallest tested extents are 2.7\,\% (inter-turn) and 2.3\,\% (inter-winding).}\label{tab:mde}
\begin{tabular}{@{}llcccc@{}}
\toprule
Fault class & Feature & PMSG MDE & PMSG DS & SCIG MDE & SCIG DS \\
\midrule
inter-turn & $V_2/V_1$ & 2.8 [2.7, 3.8] & 76 & 2.7 [2.7, 7.4] & 73 \\
inter-turn & $I_2/I_1$ & 7.4 [2.7, 7.4] & 56 & 2.8 [2.7, 7.4] & 75 \\
inter-turn & PI$_d$\,2$f_e$ & 11.6 [3.8, 11.6] & 31 & 2.7 [2.7, 7.4] & 78 \\
inter-turn & $V_q$\,2$f_e$ & 7.4 [2.8, 7.7] & 57 & 7.4 [7.4, 7.4] & 60 \\
inter-winding & $V_2/V_1$ & 2.3 [2.3, 2.3] & 94 & 2.3 [2.3, 2.3] & 100 \\
inter-winding & $I_2/I_1$ & 9.7 [2.3, 9.7] & 81 & 2.3 [2.3, 2.3] & 100 \\
inter-winding & PI$_d$\,2$f_e$ & 10.6 [3.2, 14.3] & 72 & 2.3 [2.3, 2.3] & 100 \\
inter-winding & $V_q$\,2$f_e$ & 2.3 [2.3, 9.7] & 90 & 2.3 [2.3, 9.7] & 92 \\
\botrule
\end{tabular}
\end{table}


\begin{figure}[t]
\centering
\includegraphics[width=\textwidth]{G_severity_combined}
\caption{Signal-to-drift ratio of $V_2/V_1$ and $I_2/I_1$ for the three alternatives, three-cycle windows. Left pair: against the design-level fault extent (small dots: trials; large markers: median per extent level and fault impedance, marker shape encoding the impedance). Right pair: against physical fault severity (RMS fault current relative to the RMS pre-fault stator current), all trials.}\label{fig:severity}
\end{figure}

\subsection{Operating envelope}\label{sec:res-op}

Speed dominates the operating-point dependence (Online Resource~1, Table~S3): in the PMSG the median SDR of the best feature and the detectable share rise from the lowest to the highest speed, in the SCIG less so and from a higher level; torque changes little. Two cautions apply. The dependence is consistent with a resistor-limited fault current that scales with speed, so it is a property of the emulation as much as of the machines; and the pooled null is a mixture over operating points---the share of trials whose healthy change exceeds the pooled 95th percentile ranges from 0 to 16\,\% by speed for the key features---so the 5\,\% false-alarm reading holds on average over the envelope, not at each point. Both make the lowest monitored speed a parameter of the requirement.

\subsection{Observation suites and the design decision}\label{sec:res-suites}

Table~\ref{tab:suites} and Fig.~\ref{fig:suites} give the suite results for the PMSG and SCIG at five-cycle windows, Table~S1 of Online Resource~1 for the three alternatives on the common feature set at three cycles, and Table~\ref{tab:decision} the decision table of step~5. For the PMSG, current transducers alone resolve inter-turn faults from 7.4\,\% and inter-winding faults from 9.7\,\%; three voltage transducers take these to 3.8 and 2.3\,\%, and neither controller access nor a torque transducer improves the fixed-feature result further (controller access raises the oracle share from 88 to 95\,\%). The drive-internal suite reaches 7.4\,\% for inter-turn faults through the duty-cycle negative sequence, a weaker version of the signature the measured terminal voltage carries. For the SCIG the cheapest suites already reach what any suite reaches---7.4\,\% for inter-turn faults through $I_2/I_1$ or, with no added sensor, the $d$-axis PI output, and 2.3\,\% for inter-winding faults---and additional hardware raises the detectable share rather than the reach. For the WFSG every suite detects all inter-winding faults from 2.3\,\% and 78--98\,\% of the inter-turn trials at the matching impedance, and the field current alone detects 86 and 100\,\%. The nested selection of the fixed feature changes no share by more than five points (Table~S1), so the in-sample choice is not what drives these numbers; the oracle shares, however, come with false-alarm rates of 44--78\,\% on the healthy trials for the suites with many features, so the oracle column is an upper bound only. A shaft torque transducer adds nothing for winding faults on the two machines on which it was available (the torque features' healthy relative change exceeds 100\,\% on the SCIG and their median SDR is below 0.3 on the PMSG).

\begin{table}[t]
\scriptsize
\setlength{\tabcolsep}{3pt}
\caption{Observation suites versus what they detect, PMSG and SCIG (five-cycle windows): best fixed feature of the suite, its minimum detectable extent (MDE, \%; limit-of-detection convention) and detectable share (DS, \%; per-trial oracle share in brackets). Features failing the reliability screen (three or more false alarms on the nine healthy trials, or null quantile above one) are excluded for that machine; the oracle false-alarm rates are in Table~\ref{tab:suites3}.}\label{tab:suites}
\begin{tabular}{@{}lllllll@{}}
\toprule
Suite & PMSG feature & MDE & DS [oracle] & SCIG feature & MDE & DS [oracle] \\
\midrule
\multicolumn{6}{@{}l}{\emph{Inter-turn faults}} \\
drive-internal & $D_2/D_1$ & 7.4 & 64 [87] & PI$_d$\,2$f_e$ & 7.4 & 78 [99] \\
3 CT & $I_2/I_1$ & 7.4 & 56 [77] & $I_2/I_1$ & 7.4 & 75 [87] \\
3 CT + 3 VT & $V_2/V_1$ & 3.8 & 76 [88] & $I_2/I_1$ & 7.4 & 75 [90] \\
3 CT + 3 VT + drive & $V_2/V_1$ & 3.8 & 76 [95] & PI$_d$\,2$f_e$ & 7.4 & 78 [99] \\
+ torque transducer & $V_2/V_1$ & 3.8 & 76 [95] & PI$_d$\,2$f_e$ & 7.4 & 78 [99] \\
\multicolumn{6}{@{}l}{\emph{Inter-winding faults}} \\
drive-internal & $D_2/D_1$ & 14.3 & 90 [100] & PI$_d$\,2$f_e$ & 2.3 & 100 [100] \\
3 CT & $I_2/I_1$ & 9.7 & 81 [94] & $I_2/I_1$ & 2.3 & 100 [100] \\
3 CT + 3 VT & $V_2/V_1$ & 2.3 & 94 [97] & $I_2/I_1$ & 2.3 & 100 [100] \\
3 CT + 3 VT + drive & $V_2/V_1$ & 2.3 & 94 [100] & PI$_d$\,2$f_e$ & 2.3 & 100 [100] \\
+ torque transducer & $V_2/V_1$ & 2.3 & 94 [100] & PI$_d$\,2$f_e$ & 2.3 & 100 [100] \\
\botrule
\end{tabular}
\end{table}


\begin{figure}[t]
\centering
\includegraphics[width=\textwidth]{F_sensor_suites_paper}
\caption{Observation suites versus what they detect, PMSG and SCIG (five-cycle windows): minimum detectable extent (top; the fixed feature is named) and detectable share (bottom) for the fixed best feature of each suite (filled) and the per-trial oracle (hollow).}\label{fig:suites}
\end{figure}

The decision table turns this into the selection of step~5. At $\alpha=0.95$ no suite on any alternative meets an inter-turn requirement of 3 or 5\,\% with 95\,\% bootstrap confidence; the PMSG's voltage suite is the only candidate below 5\,\% and its verdict is uncertain, a consequence of the three-cycle windows and the coarse extent grid. At 8\,\% the drive-internal suite meets the requirement on the PMSG and the SCIG, i.e.\ no added sensor is needed for either; the WFSG's inter-turn verdicts are bounded by its smallest tested extent. For inter-winding faults the SCIG and the WFSG meet every requirement from 3\,\% with the drive-internal suite, whereas the PMSG needs current transducers at 12\,\% and is uncertain below (its five-cycle result, Table~\ref{tab:suites}, reaches 2.3\,\% with voltage transducers). Relaxing $\alpha$ to 0.90 admits the PMSG's voltage suite at 5\,\% for inter-turn faults; tightening it to 0.99 removes the PMSG's current-only option at 12\,\% for inter-winding faults. The design content is therefore twofold: for a requirement around 8\,\% of a branch the three alternatives are equally diagnosable with drive-internal quantities and diagnosability does not discriminate between them; below 5\,\% only the PMSG with voltage transducers is a candidate and the evidence is not yet conclusive. In neither case would the architecture decision change on diagnosability grounds; the monitoring specification would.

\begin{table}[t]
\scriptsize
\setlength{\tabcolsep}{3pt}
\caption{Decision table: cheapest observation suite whose minimum detectable extent meets the requirement $e_{\mathrm{req}}$ with bootstrap probability $\ge0.95$, per alternative, fault class and false-alarm level $\alpha$ (three-cycle windows; suites ordered drive-internal $<$ 3\,CT $<$ 3\,CT+3\,VT $<$ 3\,CT+3\,VT+drive; exc: WFSG field current). `(suite)?': no suite meets with that confidence, the cheapest suite whose verdict is uncertain is shown; `none': every suite fails. WFSG inter-turn verdicts are bounded by its smallest tested extent at the matching impedance (11.6\,\%).}\label{tab:decision}
\begin{tabular}{@{}llcccccc@{}}
\toprule
\multirow{2}{*}{$\alpha$} & \multirow{2}{*}{$e_{\mathrm{req}}$ (\%)} & \multicolumn{2}{c}{PMSG} & \multicolumn{2}{c}{SCIG} & \multicolumn{2}{c}{WFSG} \\
\cmidrule(lr){3-4}\cmidrule(lr){5-6}\cmidrule(l){7-8}
 & & IT & IW & IT & IW & IT & IW \\
\midrule
0.90 & 3 & (3CT+3VT)? & (drive)? & none & drive & none & drive \\
0.90 & 5 & 3CT+3VT & (drive)? & none & drive & none & drive \\
0.90 & 8 & drive & (drive)? & drive & drive & none & drive \\
0.90 & 12 & drive & 3CT & drive & drive & drive & drive \\
0.95 & 3 & (3CT+3VT)? & (drive)? & none & drive & none & drive \\
0.95 & 5 & (drive)? & (drive)? & none & drive & none & drive \\
0.95 & 8 & drive & (drive)? & drive & drive & none & drive \\
0.95 & 12 & drive & 3CT & drive & drive & drive & drive \\
0.99 & 3 & none & none & none & drive & none & drive \\
0.99 & 5 & (3CT+3VT)? & none & none & drive & none & drive \\
0.99 & 8 & drive & none & drive & drive & none & drive \\
0.99 & 12 & drive & (drive)? & drive & drive & drive & drive \\
\botrule
\end{tabular}
\end{table}


\subsection{Transfer across alternatives}\label{sec:res-transfer}

Table~\ref{tab:transfer} reports step~6 for the PMSG--SCIG pair and Table~S11 of Online Resource~1 for the three alternatives; the metrics on all windows and on the out-of-reference windows agree to within 0.02, so the windows used for normalisation do not drive the result. In physical units a gradient-boosted detector reaches an AUC of 0.87--0.91 within a machine on unseen tap pairs but 0.78 (PMSG$\to$SCIG) and 0.55 (SCIG$\to$PMSG) across; with its features scaled on the target machine's healthy trials, a commissioning baseline, 0.82 and 0.75, and a logistic detector under that rule falls below chance in one direction (0.45); with each feature referenced to the first half of the trial's own reference interval, 0.91 and 0.84, against 0.91--0.94 within a machine, and a logistic detector transfers at 0.89 and 0.87. The threshold that gives 1\,\% false positives on all negatives fires on 0--3\,\% of the healthy trials' fault-time windows, the only negatives that contain the contactor operation. Between the three alternatives, self-referenced detectors transfer in all six directions with AUC 0.80--0.98 on the out-of-reference windows, and a detector trained on two machines detects the faults of the third with 0.83 (PMSG), 0.88 (SCIG) and 0.97 (WFSG); the WFSG entries are optimistic because that bench has no healthy trials and therefore no nuisance-transient negatives. Restricted to each suite's features (Online Resource~1, Table~S6), every suite transfers at 0.82--0.95 when self-referenced and at 0.62--0.90 under the commissioning baseline.

\begin{table}[t]
\scriptsize
\caption{Transfer of a window-level detector between PMSG and SCIG: AUC/true-positive rate at 1\,\% false-alarm rate. Within-machine columns use five-fold cross-validation grouped by tap pair; cross columns train on all trials of one machine and test on all trials of the other. Negatives are the reference and recovery windows of all trials and the fault-time windows of the healthy trials.}\label{tab:transfer}
\begin{tabular}{@{}llcccc@{}}
\toprule
Referencing & Model & P$\to$P & S$\to$S & P$\to$S & S$\to$P \\
\midrule
physical units & logistic & 0.80/0.56 & 0.79/0.57 & 0.71/0.41 & 0.61/0.26 \\
physical units & GBDT & 0.89/0.61 & 0.90/0.65 & 0.76/0.35 & 0.59/0.28 \\
$z$ on target's healthy trials & logistic & 0.80/0.56 & 0.79/0.57 & 0.70/0.41 & 0.39/0.07 \\
$z$ on target's healthy trials & GBDT & 0.88/0.62 & 0.90/0.66 & 0.83/0.23 & 0.75/0.31 \\
$|z|$ vs own reference & logistic & 0.91/0.71 & 0.92/0.80 & 0.91/0.70 & 0.89/0.61 \\
$|z|$ vs own reference & GBDT & 0.93/0.76 & 0.94/0.83 & 0.93/0.74 & 0.85/0.43 \\
\botrule
\end{tabular}
\end{table}


\subsection{Robustness and the two nulls}\label{sec:res-robust}

Table~S4 of Online Resource~1 shows how the tuning choices move the design quantities, with the screen applied at every setting. The inter-turn MDEs of all five key features are the same at three, five and eight cycles on both machines and the detectable shares move by at most six points; for inter-winding faults one MDE moves (the PMSG's $V_2/V_1$, 9.7\,\% at three cycles against 2.3\,\% at five and eight). Tightening the null quantile from 0.95 to 0.99 lowers the shares by 3--9 points and moves two PMSG inter-winding MDEs up one level; relaxing it to 0.90 raises the shares by up to 9 points and moves one down. No ordering of the two machines on any feature reverses, but the suite conclusion for the PMSG's inter-winding faults depends on $\alpha$ and on the window length, as the decision table shows. The screen removes between one and five features per machine and setting (Table~S5); none of the five key features is removed at any setting, and the three-cycle result of $V_2/V_1$ on the PMSG (two false alarms in nine healthy trials) passes the binomial rule. An absolute-change variant of the SDR, which removes the healthy level of the feature from numerator and null, ranks the 30 features with a Spearman correlation of 0.92--0.98 against the relative variant and picks the same three top features in every machine and class. The bootstrap intervals (Table~\ref{tab:mde}; Online Resource~1, Fig.~S2) show the tap-pair-level variability behind each number: the median SDRs are known to within a factor of two to four, the shares to $\pm$8--10 points, and the ratios of median SDR exclude one where the signature relocates and include it where it does not.

Table~S8 of Online Resource~1 quantifies the assumption that the healthy reference lies immediately before the fault. When the reference is instead the healthy trial recorded at the same operating point at another time---a commissioning baseline---the null widens by a factor of two to seven, the median SDR of the inter-turn features falls to about one (PMSG $V_2/V_1$ from 3.0 to 1.1; SCIG $I_2/I_1$ from 4.5 to 1.1) and their detectable shares to 44--59\,\%, while inter-winding faults remain detectable in 67--94\,\% of trials. The MDEs reported here are therefore statements about a short-horizon change detector; a monitor that compares against an old baseline needs the wider null, and the transfer result says the same thing from the other side.
